\documentclass[10pt,a4paper]{amsart}
\usepackage[margin=19mm]{geometry}
\usepackage{amsmath,amssymb,mathtools}
\usepackage[scr=rsfs,cal=euler]{mathalfa}
\usepackage{xfrac}
\usepackage[unicode,hidelinks,bookmarks=false]{hyperref}
\newcommand{\ie}{i.e.\ }
\newcommand{\cf}{cf.\ }
\newcommand{\resp}{respectively\ }
\newcommand{\Subsection}{\subsection}
\providecommand{\nobreakdash}{\nobreak\hbox{-}}
\setlength{\emergencystretch}{2em}
\multlinegap0pt
\usepackage{tikz}
\usepackage{mathtools}
\usepackage{enumerate}
\usepackage{amssymb}
\usepackage{bm}
\usepackage{xcolor}
\usepackage{stackengine}
\newtheorem{lemma}{Lemma}
\newcommand{\CS}{\mathcal{S}}
\newcommand{\avec}[2]{\norm{#1}_{#2,P}}
\newcommand{\bfs}{\mathbf{s}}
 \newcommand{\norm}[1]{\lVert#1\rVert}
\title{Query Learning Nearly Pauli Sparse Unitaries in Diamond Distance}
\author{Zahra Honjani, Mohsen Heidari}
\date{Source: arXiv:2604.00203v1 (CC BY 4.0)}
\begin{document}
\maketitle

\noindent\emph{Source and licence.} Query Learning Nearly Pauli Sparse Unitaries in Diamond Distance. Version arXiv:2604.00203v1 (CC BY 4.0), distributed by arXiv under the Creative Commons Attribution 4.0 International licence, http://creativecommons.org/licenses/by/4.0/, as recorded in the arXiv licence metadata for this exact version and shown on the abstract page https://arxiv.org/abs/2604.00203v1. Full licence notice: \texttt{licenses/arxiv-2604.00203-CC-BY-4.0.txt}.\par\smallskip

\noindent\textit{Editorial scope.} Counted unit: the article's lemma lm:app\_unitary (source0.tex lines 2543--2561). The article's complete original proof of it is delivered with it (source0.tex lines 2563--2657, one \texttt{proof} environment of 95 lines). No mathematics is written, repaired or re-derived: the statement and the proof are the article's own lines, in the article's own notation, and no retained line is substituted or re-flowed.  No further source-local context is needed: every cross-reference of every retained line resolves inside the two delivered spans.  Nothing was omitted from any retained span, so the package carries no omission record.  No retained line contains a citation, so the wrapper carries no bibliography and no bibliography entry.  No retained line contains a figure environment, an image-inclusion command or drawing code, so no image is delivered and the package declares no asset dependency.  Editorial formatting only, in addition: an AMS \texttt{amsart} preamble that restates the article's own declarations and macros used by the retained lines, namely the environments \texttt{lemma} on separate counters, together with the standard packages those lines need.  The retained environments print in this document as Lemma 1 in order of appearance, whereas the article prints the same items with the numbering of its own full document.  The complete native source of the article as distributed in the arXiv e-print for this exact version is bundled unchanged as \texttt{sources/197574/source0.tex}.
    \begin{lemma}
        Suppose an $n$-qubit unitary $U$ is approximated by $\tilde{U}$ such that 
        
        
        $\norm*{U - \tilde{U}}_{2}^2 \le \varepsilon/4$ and 
        $\norm*{\tilde{U}}_{1,P} = L_1$. 
        Then, $U$ can be approximated in $\ell_2$-norm by a $\tfrac{4L_1^2}{\varepsilon}$-sparse operator $\hat{U}$ generated from the Pauli coefficients $U$ larger than $\tfrac{\varepsilon}{2L_1}$. 
         More precisely, let $$S = \{\bfs\in \{0,1,2,3\}^n \mid |\alpha_{\bfs}| \ge {\frac{\varepsilon}{2L_1}} \}.$$ 
         Then, there exists a set $\mathcal{S}^* \subseteq S$,  with size $|\mathcal{S}^*| \le \tfrac{4L_1^2}{\varepsilon}$, such that the operator 
        \[
            \hat{U} = \sum_{\bfs \in \mathcal{S}^*} \hat{\alpha}_{\bfs} \sigma^{\bfs},
        \]
        with  $|\hat{\alpha}_{\bfs} - \alpha_{\bfs}| \le \frac{\varepsilon}{2L_1}$ for all $\bfs \in \CS^*$, 
        satisfies 
        \[
            \avec{U - \hat{U}}{2}^2 \le 6\varepsilon.
        \]
        \label{lm:app_unitary}
    \end{lemma}

    \begin{proof}
        We first show existence of a sparse $V$ that approximates $\tilde{U}$ and hence $U$.
        Write the Pauli expansion of $\tilde{U} = \sum_{\bfs} \tilde{\alpha}_{\bfs} \sigma^{\bfs}$. 
        Let 
        $S' = \{\bfs \mid |\tilde{\alpha}_{\bfs}| \ge \frac{\varepsilon}{4L_1}\}$ 
        and define 
        $V = \sum_{\bfs \in S'} \tilde{\alpha}_{\bfs} \sigma^{\bfs}$. 
        Then $$\avec{V}{1}\le \norm*{\tilde{U}}_{1,P} = L_1$$ and 
        $|S'| \le \frac{4L_1^2}{\varepsilon}$.
        Moreover,
        \[
        \avec{\tilde{U} - V}{2}^2
        = \sum_{\bfs \notin S'} |\tilde{\alpha}_{\bfs}|^2
        \le \max_{\bfs \notin S'} |\tilde{\alpha}_{\bfs}| 
            \sum_{\bfs \notin S'} |\tilde{\alpha}_{\bfs}| 
        \le \frac{\varepsilon}{4L_1} L_1 = \frac{\varepsilon}{4}.
        \]
        Using $\|A+B\|^2 \le 2(\|A\|^2 + \|B\|^2)$, we get
        \[
        \avec{U - V}{2}^2
        \le 2(\avec{U - \tilde{U}}{2}^2 +\avec{\tilde{U} - V}{2}^2)
        \le \varepsilon.
        \]


        It remains to show that the sparse approximation $\hat{U}$ constructed from $U$ satisfies the desired error bound.
        
        
        Particularly, we show that  
        $\norm*{U - \hat{U}}_{2,P}^2 \le 3\varepsilon$.
        
        From the above construction, we have
        \[
        \varepsilon \ge \avec{U - V}{2}^2 
        = \sum_{\bfs \notin S'} |\alpha_{\bfs}|^2 
          + \sum_{\bfs \in S'} |\alpha_{\bfs} - \tilde{\alpha}_{\bfs}|^2.
        \]
        implying $\sum_{\bfs \notin S'} |\alpha_{\bfs}|^2 \le \varepsilon$. 
        Let $\CS^* = S \cap S'$ and 
        $K = \sum_{\bfs \in \CS^*} \alpha_{\bfs} \sigma^{\bfs}$. 
        Then
        
        \begin{align*}
            \avec{U - K}{2}^2 &=\sum_{\bfs \notin \CS^*} |\alpha_{\bfs}|^2\\
            &
            = \sum_{\bfs \notin S'} |\alpha_{\bfs}|^2 
              + \sum_{\bfs \in S' - \CS^*} |\alpha_{\bfs}|^2 \\
            &\le \varepsilon + |S'| \cdot \frac{\varepsilon^2}{4L_1^2} \le 2\varepsilon.
        \end{align*}        
        Since $\CS^* \subseteq S',$ $|\CS^*|\leq \tfrac{4L_1^2}{\varepsilon}$. Moreover, as $\gamma \leq \tfrac{\varepsilon}{2L_1}$, the error of approximating $U$ by $\hat{U}$ can be bounded as follows.   
        \begin{align*}
            \avec{U - \hat{U}}{2}^2 \leq 2(\avec{U - K}{2}^2 + \avec{K - \hat{U}}{2}^2)\leq 2(2\varepsilon + |\CS^*| \gamma^2) \leq 6\varepsilon.
        \end{align*}


        
        
        
        
        
        
        
        
        
        
        
        
        
        
        
        
        
        
        
        
        
        
        
        
        
        
        
        
        
        
        
        
        
        
        
        
        
        
        
    \end{proof}

\end{document}
